% Abstract: one page, unnumbered, listed in the contents.
\addchap{Abstract}

Stationary lithium-ion batteries now balance a growing share of the
electricity grid, yet the operators who own them rarely know how much
of their capacity remains. Laboratory ageing models are calibrated on
small cells under conditions that a grid container never sees. This thesis asks whether the state of health of grid-scale storage
can be learned directly from the operational data that every site already
records.

Part~I assembles the evidence. Three years of telemetry from eleven
installations across the Low Countries, totalling 214\,MWh of nameplate
capacity, are aligned with periodic reference capacity tests. Their
analysis shows that calendar ageing at high state of charge, not cycling
depth, explains most of the observed fade, and that the commonly used
equivalent full cycle is a poor proxy for degradation under
frequency-response duty.

Part~II develops a hybrid estimator that embeds a semi-empirical ageing
law inside a Gaussian-process correction learned from partial charge
segments. On held-out sites it estimates capacity to within 0.9\,per cent
root-mean-square error, and it reports a calibrated uncertainty that widens when a site
leaves the operating envelope of its training data. A field trial at a
12\,MWh installation shows that the estimate can safely relax a
conservative state-of-charge window, recovering 6.4\,per cent of usable
energy without measurable additional ageing over nine months.

\bigskip\noindent
{\sansbold\small\color{accent}\spacedcaps{Keywords}}\enspace
{\sffamily\small lithium-ion batteries, state of health, grid storage,
Gaussian processes, calendar ageing, frequency response}
